\section*{الفصل \ref{ch:b2:phanerozoic-earth} --- الأرض في دهر الحياة الظاهرة}

\begin{solution}{exo:b2:phanerozoic-earth:1}
(أ) باريس (\ang{5.1} شمالًا) ونيويورك (\ang{3.6} جنوبًا)؛ أمّا كيب تاون وسيدني فكانتا قرب \ang{60} جنوبًا.
(ب) كيب تاون: $61.9 - 33.9 = \ang{28.0}$، و$28.0 \times 111.2 = \qty{3114}{km}$؛ ونيويورك:
$40.7 + 3.6 = \ang{44.3}$، أي \qty{4926}{km}. (ج) $3114/300 = \qty{10.4}{km}$ لكل مليون سنة،
أي \qty{1.0}{cm} في السنة؛ و$4926/300 = \qty{16.4}{km}$ لكل مليون سنة، أي \qty{1.6}{cm} في السنة.
\end{solution}

\begin{solution}{exo:b2:phanerozoic-earth:2}
(أ) 72 و120.1 و50.5 و\qty{135.4}{Myr}؛ والمتوسط $(444 - 66)/4 = \qty{94.5}{Myr}$. (ب) نهاية
البرمي (سيبيريا، بفارق \qty{1.9}{Myr})، ونهاية الترياسي (CAMP، \qty{0.4}{Myr}) ونهاية
الطباشيري (الدكن، \qty{0.5}{Myr}). (ج) لا: فالفترات تتراوح من 50 إلى
\qty{135}{Myr}؛ والمتوسط موجود لأيّ مجموعة من التواريخ.
\end{solution}

\begin{solution}{exo:b2:phanerozoic-earth:3}
(أ) $\qty{91}{m}/\qty{0.036}{mm.yr^{-1}} = \num{2.5e6}$ سنة. (ب) $3.7/0.036 = 103$ مرة. (ج)
بعيدًا عن أيّ شاطئ، في ماء صافٍ، كان القليل من طين الأنهار يبلغ القاع؛ وكان مطر
أصداف الطحالب هو الراسب الرئيسي.
\end{solution}

\begin{solution}{exo:b2:phanerozoic-earth:4}
(أ) من الحجر الرملي النهري إلى الطباشير تتعمّق الرواسب نحو الأعلى: إنه تجاوز بحري.
ورمل الشاطئ فوق الطباشير يسجّل انحسارًا. (ب) لا: فهو يسجّل ارتفاعًا في مستوى البحر النسبي،
قد ينتجه أيضًا هبوط المكان، أو نقص في إمداد الرواسب. (ج)
العثور على التجاوز نفسه، في الزمن نفسه، في قارّات أخرى، أو إزالة
الهبوط \omterm{def:b2:basins-compared:backstripping}{بالتجريد العكسي}.
\end{solution}

\begin{solution}{exo:b2:phanerozoic-earth:5}
(أ) $233/97 = \qty{2.4}{m}$ لكل مليون سنة، أي \qty{2.4}{\micro m} في السنة. (ب) $3.7/0.0024 =
1540$: فارتفاع اليوم أسرع بنحو 1500 مرة. (ج) المنحنى المستخرج \omterm{def:b2:basins-compared:backstripping}{بالتجريد العكسي}: فقد أظهر الحفر
أن مخطّطات الدورات تصيب في توقيت الانخفاضات لكنها تبالغ في حجمها، بعامل اثنين أو
أكثر؛ ولنيوجيرسي انحيازها الخاصّ (التضاريس الديناميكية)، ومن هنا التقدير $100 \pm
\qty{50}{m}$.
\end{solution}

\begin{solution}{exo:b2:phanerozoic-earth:6}
(أ) $\num{e5} \times \num{e9} \times 2900 = \qty{2.9e17}{kg}$ من البازلت، و$\num{2.9e17}
\times 0.0075 = \qty{2.2e15}{kg}$ من ثاني أكسيد الكربون، و$\num{2.2e15}/0.04401 = \qty{4.9e16}{mol}$.
(ب) $\num{278e-6} \times (44.01/28.97) \times \num{5.1e18} = \qty{2.15e15}{kg}$. (ج) أطلقت نبضة
واحدة قدر ما حمله الهواء كلّه عام 1750 تقريبًا. (د) $\num{e17}/\num{2.15e15} = 46.5$: قرابة خمسين.
\end{solution}

\begin{solution}{exo:b2:phanerozoic-earth:7}
(أ) $2630 \times \frac{\pi}{6} \times (\num{e4})^3 = \qty{1.38e15}{kg}$. (ب) $\tfrac12 \times
\num{1.38e15} \times (\num{11.2e3})^2 = \qty{8.7e22}{J}$؛ وعند \qty{20}{km.s^{-1}}،
\qty{2.76e23}{J}. (ج) سنة من \qty{47}{TW} هي $\num{47e12} \times \num{3.156e7} =
\qty{1.48e21}{J}$: أي 59 و186 سنة.
\end{solution}

\begin{solution}{exo:b2:phanerozoic-earth:8}
(أ) اللون الداكن، والتصفّح الدقيق، وغياب الجحور، ووفرة المادة العضوية، والبيريت؛ وكثيرًا ما يرافقها انزياح في
نظائر الكربون وعلامات احترار. (ب) حدث التوارسي مع كارو--فيرار، و\babelsublr{OAE~1a} مع
أونتونغ جاوة. (ج) تحمل التواريخ شكوكًا تبلغ ملايين السنين، والأحداث والمقاطعات
كثيرة، فالتزامنات العرضية مرجّحة؛ والرابط السببي يحتاج إلى أن تبدأ المقاطعة
قبيل الحدث وإلى متتبّع للنشاط البركاني، كالزئبق، في الرواسب نفسها.
\end{solution}

\begin{solution}{exo:b2:phanerozoic-earth:9}
(أ) $60.9 + 5.1 = \ang{66.0}$، و$66.0 \times 111.2 = \qty{7339}{km}$ في \qty{200}{Myr}:
أي \qty{36.7}{km} لكل مليون سنة، و\qty{3.7}{cm} في السنة. (ب) أكثر من ضعف
\qty{1.6}{cm} في السنة منذ \qty{300}{Ma}. (ج) لا يُعرف إلا تغيّر خطّ العرض؛ أمّا تغيّر
خطّ الطول، الحرّ في إطار مغناطيسي قديم، وأيّ التفاف فيضيفان إلى المسافة.
\end{solution}

\begin{solution}{exo:b2:phanerozoic-earth:10}
(أ) كسر القاع الذي عمره $t$ هو $p(t) = \frac2T(1 - t/T)$. ومن ثَمّ $\int_0^T tp\,\mathrm{d}t =
\frac2T\bigl(\frac{T^2}2 - \frac{T^2}3\bigr) = \frac T3$ و$\int_0^T \sqrt t\,p\,\mathrm{d}t =
\frac2T\bigl(\frac23 T^{3/2} - \frac25 T^{3/2}\bigr) = \frac8{15}\sqrt T$. (ب) $T = 3 \times 64.2 =
\qty{192.6}{Myr}$؛ و$2500 + \frac8{15} \times 350 \times 13.878 = \qty{5091}{m}$. (ج) $T/2 =
\qty{96.3}{Myr}$: \qty{4332}{m}؛ والانخفاض \qty{759}{m}، أي أقلّ بنسبة \qty{2}{\percent} من
\qty{774}{m}. (د) النموذج المنتظم: $\frac23\sqrt\tau = \frac23\sqrt{2\bar t} = 0.943\sqrt{\bar t}$؛
والنموذج المثلّثي: $\frac8{15}\sqrt{3\bar t} = 0.924\sqrt{\bar t}$. فالجواب يتوقّف على العمر المتوسط أكثر
بكثير مما يتوقّف على الشكل.
\end{solution}

\begin{solution}{exo:b2:phanerozoic-earth:11}
(أ) في $\Delta s$، العامل $1 - f^{-1/2}$ وحده يتوقّف على $f$؛ فاقسم على قيمته عند $f = 2$،
$1 - 2^{-1/2} = 0.2929$. (ب) $f = (1 - 0.2929\,\Delta s/432)^{-2}$: 1.071 من أجل \qty{50}{m}، و1.239
من أجل \qty{150}{m}؛ أي أسرع من اليوم بنسبة 7 إلى \qty{24}{\percent}. (ج) $432 \times (1 -
0.8^{-1/2})/0.2929 = \qty{-174}{m}$: فينخفض البحر بمقدار \qty{174}{m}. (د) لا يتغيّر توزيع
الأعمار إلا كلّما اندسّ القاع القديم واستُبدل، وهذا يستغرق زمنًا من رتبة $\tau$،
نحو \qty{100}{Myr}؛ فيتبعه البحر على مدى عشرات ملايين السنين.
\end{solution}

\begin{solution}{exo:b2:phanerozoic-earth:12}
(أ) $2^{-7.4/2} = 0.077$: أي زال \qty{92}{\percent}. (ب) $t = 2\log_2(1/0.19) = \qty{4.8}{Myr}$.
(ج) $4.8/0.060 = 80$ مرة أسرع. (د) في طابق طويل يزيل انقراض الخلفية وحده
معظم الأنواع، فتُضاف نبضة قصيرة إلى مجموع كبير ويصعب فصلها عنه؛ كما أن
نقص السجلّ يزحزح آخر الظهورات نحو الأقدم.
\end{solution}

\begin{solution}{pb:b2:phanerozoic-earth:1}
\textbf{1.} 2500 و4250 و5304 و\qty{6000}{m}.
\textbf{2.} $\tau = 2 \times 64.2 = \qty{128.4}{Myr}$؛ و$2500 + 350\sqrt{128.4} = \qty{6466}{m}$.
\textbf{3.} $\num{293e6}/(65000 \times 128.4) = \qty{35.1}{km}$ لكل مليون سنة،
أي \qty{35.1}{mm} في السنة: معدّل بطيء.
\textbf{4.} $\frac1\tau\int_0^\tau(2500 + 350\sqrt t)\,\mathrm{d}t = 2500 + \frac23 \times
350\sqrt\tau = \qty{5144}{m}$.
\textbf{5.} قانون الجذر التربيعي أعمق مما ينبغي للقاع الأقدم من بضع عشرات ملايين السنين،
إذ يستوي هذا القاع؛ والأعمار الحقيقية ليست موزّعة بانتظام، فالقاع الفتيّ هو الأوفر؛
والرواسب والجبال البحرية والهضاب ترفع القاع.
\textbf{6.} $\num{293e6} \times 5.144 = \qty{1.51e9}{km^3}$، أي أكثر بنسبة \qty{13}{\percent} من
المحيط كلّه: فالنموذج أعمق مما ينبغي.
\textbf{7.} $v = \qty{17.6}{km}$ لكل مليون سنة؛ و$17.6 \times 128.4 = \qty{2260}{km}$ على
كلّ جانب.
\textbf{8.} $t = 1000/17.6 = \qty{56.8}{Myr}$؛ و$2500 + 350\sqrt{56.8} = \qty{5138}{m}$.
\textbf{9.} على جانب واحد، مع $x = vt$: $\int_0^{v\tau}\bigl(d(\tau) - d(x/v)\bigr)\mathrm{d}x
= vc\int_0^\tau(\sqrt\tau - \sqrt t)\,\mathrm{d}t = \frac13 vc\,\tau^{3/2}$؛ وضعف ذلك للجانبين.
$\frac23 \times 17.6 \times 0.350 \times 128.4^{3/2} = \frac23 \times 17.6 \times 0.350 \times
1454.9 = \qty{5975}{km^2}$.
\textbf{10.} $65000 \times 5975 = \qty{3.88e8}{km^3}$.
\textbf{11.} $\num{3.88e8}/\num{1.33e9} = 0.29$.
\textbf{12.} قاعها أفتى في المتوسط، فهو أسخن وأقلّ كثافة وأعلى؛ وشريط القاع الفتيّ
على امتداد المحور أعرض.
\textbf{13.} $\tau = \qty{64.2}{Myr}$؛ و$2500 + \frac23 \times 350\sqrt{64.2} = \qty{4370}{m}$.
\textbf{14.} $5144 - 4370 = \qty{774}{m}$.
\textbf{15.} $\num{293e6} \times 0.774 = \qty{2.27e8}{km^3}$.
\textbf{16.} $\num{2.27e8}/\num{362e6} = \qty{0.627}{km}$، أي \qty{627}{m}.
\textbf{17.} $\Delta s = 627 \times (3300 - 1025)/3300 = \qty{432}{m}$.
\textbf{18.} $f = \bigl(1 - 0.2929 \times 100/432\bigr)^{-2} = 1.15$: أي أسرع بنسبة \qty{15}{\percent}.
\textbf{19.} $0.149 \times 148 = \qty{22.1}{\mega km^2}$.
\textbf{20.} $48.1 + 0.32 \times 7.7 = \qty{50.6}{\percent}$، أي \qty{74.9}{\mega km^2}.
\textbf{21.} كلّما ارتفع غطّى مساحة أكبر، فيصنع الحجم نفسه طبقة أرقّ
(\cref{prop:b2:phanerozoic-earth:flooding}).
\textbf{22.} المساحات المغمورة عند 100 و200 و300 و\qty{400}{m}: 22.1 و41.7 و57.3
و\qty{71.2}{\mega km^2}. والخطوات: $(362 + 11.0) \times 0.1 = \qty{37.3e6}{km^3}$، ثم 39.4 و41.2
و\qty{42.6e6}{km^3}؛ والمجاميع التراكمية 37.3 و76.7 و117.9 و\qty{160.5e6}{km^3}. والحجم الذي يجب وضعه
\qty{156.4e6}{km^3}، بين 300 و\qty{400}{m}: $300 + (156.4 - 117.9)/42.6 \times 100
= \qty{390}{m}$.
\textbf{23.} $38.7 + 0.90 \times 9.4 = \qty{47.2}{\percent}$ من يابسة اليوم.
\textbf{24.} مضاعفة معدّل الانتشار ترفع البحر بنحو \textbf{\qty{390}{m}}
وتُغرق قرابة نصف يابسة اليوم.
\end{solution}
